\chap{开放问题与未来方向}
\label{sec:outlook}
\subsection{长期方向与组织适应}
长期智能开发仍面临目标漂移问题。人可能改变偏好，团队可能把局部指标当成最终目标，历史经验也可能固化已经失效的假设。系统需要能识别重要分歧，并在合适情境中让人重新判断。交互投影可以降低这一过程的成本，但其选择机制本身也需要可检查，避免只呈现容易解释或已经成功的部分。

组织适应同样是开放问题。成员能力增强后，合理职责范围可能扩大；项目复杂度增加后，又可能需要更多专业协作。如何根据真实工作调整团队，同时保持知识连续性和接口稳定，需要长期观察。临时成员与持续成员之间的边界，也可能随任务和模型条件变化。

经验维护需要处理可信度与遗忘。系统应保留有用的条件化知识，能够在反例出现时修订，并降低过时材料进入当前上下文的概率。记忆容量增加不能直接解决这一问题，关键仍是选择、核对和使用过程。未来可以研究经验在不同模型和团队之间迁移时的适用条件。

\subsection{生成交互的稳定性与个体差异}
生成界面的有效形式难以穷举。高层原则与典型指南提供起点，具体效果需要在人类使用中形成。持续变化的界面可能改善特定任务，却增加定位和学习成本。如何保持稳定的对象语义与常用入口，同时允许局部表达变化，是值得深入研究的设计问题。

不同人的背景和参与方式也会影响表达。架构负责人、产品使用者和临时审阅者可能关注不同问题。系统可以根据当前任务和明确偏好调整视图，但需要避免将过去一次选择当成永久规则。表达经验应能够被人纠正，并在共享协作中保持可交流的基础。

\subsection{运行模块与开发成员的进一步结合}
更远的方向是，软件运行模块本身具有智能能力，并参与自身或相关模块的开发。一个运行中的任务管理 Agent 可能持续观察使用问题，提出协议改进，并与开发成员协作。这会使运行反馈与开发知识更加接近，也会增加职责与变更边界的复杂性。

研究这种结合时，需要分别理解当前运行承诺和未来修改提案。运行模块提出的建议可以进入开发与验证过程，实际变化仍需依据明确目标和证据实施。自我修改是否带来收益，要看它能否改善真实任务、保持可理解的状态，并避免局部优化损害整体目标。本文将其作为远期研究方向。

\subsection{结语}
智能软件开发的长期目标，是让不断增强的 AI 能力通过组织、积累、反馈和交互转化为更大的工程能力。软件结构决定局部工作能否有效组合，持续成员让经验跨任务延续，局部交付和自适应验证使判断接触现实，而智能交互帮助人以较少投入理解和控制系统的发展。

这套框架中的具体机制可以变化。更强模型可能减少某些交接，更好的工具可能降低验证成本，新的表达方式可能改善人的参与。持续有效的评价问题是：系统是否更准确地理解目标，是否更自主地完成有质量的工作，以及人是否更容易发现偏差并作出有效引导。围绕这些问题开展实现与经验研究，可以逐步形成随 AI 能力共同发展的软件工程方法。
